% !TEX root = chapter-standalone.tex
\section{Counting orders}

\linkvideo{spring2021-tradeoffs:tradeoffs:orders:counting-orders} % Counting orders

Let's count the number of distinct posets on a finite set, ``up to isomorphism''.
This means we consider two posets the same if they have the same structure, regardless of the specific labels of the elements.
For small numbers of elements, we have:
\begin{itemize}
    \item For a 1-element set, there is only 1 poset (\cref{fig:singleton}), because there is only one way to order a single element.
    \item For a 2-element set, there are 2 posets (\cref{fig:twoelementspos}), because there are 2 possible orderings of a 2-element set.
    \item For a 3-element set, there are 5 posets (\cref{fig:threeelementspos}).
    \item For a 4-element set, there are 16 posets (\cref{fig:fourelementspos}).
\end{itemize}
These are illustrated in \cref{fig:posets-up-to-four}.

\vfill

\begin{figure}[h]
    % \hfill
    \subfloat[The singleton poset.
        \label{fig:singleton}
    ]{
        \rule{2cm}{0pt}\middlesag{40_dpcatfig_singleton}\rule{2cm}{0pt}
    }
    % \hfill
    \subfloat[All posets on 2-element sets, up to isomorphism.
        \label{fig:twoelementspos}
    ]{
        \centering
        \setlength{\tabcolsep}{20pt}
        \rule{1cm}{0pt}
        \begin{tabular}{cc}
            \middlesag{70_pos_2_1} & \middlesag{70_pos_2_2}
        \end{tabular}
        \rule{1cm}{0pt}
    }

    % \medskip
    \subfloat[All posets on 3-element sets, up to isomorphism.
        \label{fig:threeelementspos}]{
        \setlength{\tabcolsep}{20pt}
        \begin{tabular}{ccccc}
            \middlesag{70_pos_3_1} &
            \middlesag{70_pos_3_2} &
            \middlesag{70_pos_3_3} &
            \middlesag{70_pos_3_4} &
            \middlesag{70_pos_3_5}
        \end{tabular}
    }

    % \medskip
    \subfloat[All posets on a 4-element set, up to isomorphism.
        \label{fig:fourelementspos}
    ]{
        \centering
        \setlength{\tabcolsep}{6pt}
        \maxsizebox{\marginparwidth+\textwidth}{!}{
            \begin{tabular}{cccccc}
                \middlesag{70_pos_1}
                 & \middlesag{70_pos_2}
                 & \middlesag{70_pos_3}
                 & \middlesag{70_pos_4}
                 & \middlesag{70_pos_5} \\[+30pt]
                \middlesag{70_pos_6}
                 & \middlesag{70_pos_7}
                 & \middlesag{70_pos_8}
                 & \middlesag{70_pos_14}
                 & \middlesag{70_pos_15} \\[+30pt]
                \middlesag{70_pos_11}
                 & \middlesag{70_pos_12}
                 & \middlesag{70_pos_13}
                 & \middlesag{70_pos_9}
                 & \middlesag{70_pos_10}
                 & \middlesag{70_pos_16}
            \end{tabular}
        }
    }
    \caption{Examples of posets on sets with 1, 2, 3, and 4 elements, up to isomorphism.
        The subfigures show: (a) the single poset on a 1-element set, (b) the 2 posets on a 2-element set, (c) the 5 posets on a 3-element set, and (d) the 16 posets on a 4-element set.
    }
    \label{fig:posets-up-to-four}
\end{figure}
